\section{Materials and methods}
\label{sec:methods}

\subsection{Dataset and splits}

CAMUS~\cite{leclerc2019camus} contains apical two-chamber (2CH) and
four-chamber (4CH) views for 500 patients, annotated at ED and ES for three
structures: LV endocardium ($\ell{=}1$), LV myocardium ($\ell{=}2$) and left
atrium ($\ell{=}3$). We use the official 400/50/50 patient-level split. We
verified it byte-identical to the \texttt{subgroup\_training},
\texttt{subgroup\_validation} and \texttt{subgroup\_testing} lists released with
the dataset, so results here are directly comparable with published work that
uses the official protocol. This yields $400\times2\times2 = 1600$ training and
$200$ test frames.

Each patient additionally carries a fully annotated half cardiac cycle. Exported
at every annotated frame this gives $17{,}006$ training images, used in
Sec.~\ref{sec:results} as a data-scaling factor. We confirmed that no test or
validation patient appears in that expansion.

Images are exported at native resolution with per-image min--max normalisation.
No resampling to a common grid occurs at any stage of dataset construction, so
the pixel spacing recorded in the NIfTI header --- typically $0.308$\,mm ---
remains valid for boundary metrics.

\subsection{The annulus problem}
\label{sec:encoding}

The myocardium is the ring between the endocardial and epicardial contours: an
annulus, and therefore not representable as the single closed polygon that an
instance-segmentation label must be. Two encodings exist.

\paragraph{Filled-contour decomposition} We define the second trained class as
the \emph{epicardial disc}, $\ell \in \{1,2\}$, which is simply connected. All
three trained classes are then blob-shaped and the mask head never has to
represent a hole. At inference the CAMUS label map is recomposed by painting the
predicted disc as myocardium and overwriting its interior with the predicted
endocardium,
\begin{equation}
\hat{M}_{\text{myo}} \;=\; \hat{M}_{\text{epi}} \setminus \hat{M}_{\text{endo}},
\label{eq:decomp}
\end{equation}
so the hole emerges from the interaction of two predictions rather than from a
single mask.

\paragraph{Seam-cut annulus} Alternatively the ring is written as one polygon
that traverses the outer contour, crosses to the inner contour along a
zero-width seam, traverses it in the opposite winding, and returns. Scan-line
filling honours this and yields a genuine annulus; we verified that the
rasterisation routine used to build training targets reproduces the ring at
$0.994$ Dice, and not the filled disc ($0.49$--$0.79$).

We note that a third apparent option --- taking only the \emph{outer} contour of
the ring --- is not a distinct encoding. The outer boundary of the myocardial
ring \emph{is} the epicardial contour, making it identical to the filled
decomposition.

\subsection{Encoding ceiling}
\label{sec:ceiling}

Both encodings are lossy at a finite vertex budget, so we measure their ceiling
directly rather than assume it is negligible. Ground-truth label maps are
converted to polygons exactly as the training targets are, rasterised back, and
scored against the untouched native ground truth. This bounds any model trained
on that representation. We report two quantities: the Dice of the recomposed
four-class map, which is what evaluation ultimately sees, and the Dice of the
raw second-class target, which is what the mask head is asked to reproduce. The
second is the discriminating one, because Eq.~\ref{eq:decomp} re-carves the hole
under either encoding and therefore conceals the difference.

\subsection{Evaluation protocol}
\label{sec:protocol}

To keep our numbers comparable with the baselines evaluated on this split we
reproduce their protocol rather than improve on it:

\begin{itemize}
\item Overlap metrics (Dice, IoU) are computed per frame and per class on a
      $256\times256$ grid with a smoothing constant of $10^{-6}$.
\item Boundary metrics (HD95, ASSD) are computed at \emph{native} resolution
      using each patient's true pixel spacing, and are reported in millimetres.
      Computing them on a resampled grid while retaining the native spacing
      under-reports distances by the resampling factor.
\item Predicted instance masks are resampled to native resolution by nearest
      neighbour before scoring; for each class the single highest-confidence
      instance is kept.
\end{itemize}

We verified this reimplementation against the reference metric code on perturbed
ground truth: HD95 agrees to $2\times10^{-12}$\,mm and ASSD exactly. Our
implementation streams the pairwise boundary-distance computation in chunks
rather than materialising the full distance matrix, which keeps
native-resolution evaluation within the memory budget of the same laptop.

\subsection{Latency protocol}
\label{sec:latency-protocol}

Published latencies for the baselines were measured on other hardware and cannot
be mixed with ours. We therefore re-time every model under a single protocol:
batch size 1, CUDA-synchronised around each iteration, ten warm-up iterations
discarded, median of fifty timed iterations reported (the median rather than the
mean, because scheduler noise on a laptop skews the latter).

Three quantities are separated, because input resolution, architecture and
framework overhead are otherwise confounded. \emph{Raw forward} is the bare
module forward pass with no pre- or post-processing, measured for every model at
both $256$ and $640$\,px so that resolution is isolated from architecture; this
requires no retraining and no accuracy is claimed at the off-design size.
\emph{End-to-end} is the full inference call including letterboxing and
native-resolution mask assembly. \emph{ONNX} is the exported graph under ONNX
Runtime's CUDA provider.

\subsection{Ejection fraction}
\label{sec:ef-method}

EF is defined across the pair of cardiac phases,
$\text{EF} = (\text{EDV}-\text{ESV})/\text{EDV}$, and therefore requires pairing
each patient's ED and ES predictions in both views. Volumes are computed with
Simpson's biplane method of disks using the official CAMUS implementation, at
native resolution with each view's true spacing. We verified our copy of that
code bit-identical to the reference notebook published with the dataset: across
all 50 test patients the maximum absolute difference in both EDV and EF is
$0.0000$.

We additionally compute an \emph{oracle}: EF derived from the ground-truth masks
through the identical pipeline. This separates the error contributed by
segmentation from the irreducible gap between mask-derived EF and the clinically
recorded value, and it is the reference point against which model EF should be
read.

\subsection{Training configuration}

Unless stated otherwise we train YOLO26n-seg initialised from COCO weights at
$640$\,px, batch 8, for 100 epochs with early stopping (patience 30), on a
single NVIDIA RTX~4060 Laptop GPU (8\,GB). YOLO26 is the newest generation with
published segmentation weights and is end-to-end: it predicts a fixed instance
set with no non-maximum suppression, which removes a data-dependent stage from
the latency budget.

We use the framework's automatic optimiser selection, which for this workload
chooses MuSGD at a learning rate of $0.01$. We note explicitly that this setting
\emph{overrides} any user-specified learning rate; learning rate is therefore
varied only in the dedicated arm of Sec.~\ref{sec:results}, and readers
reproducing this work should be aware that the corresponding flag is otherwise
inert.

Several components of the stock augmentation recipe have no anatomical
justification on sector ultrasound: mosaic composites sectors into images that
cannot physically occur, horizontal flipping violates the left/right convention
of an apical view, and hue and saturation jitter are meaningless on grayscale
input. We therefore adopt an echo-aware default --- mosaic, mixup, copy-paste
and flips disabled; mild rotation ($\pm10^\circ$), translation, scaling and
brightness jitter retained. We stress that this recipe is chosen on principle
and on cost, not because it measures better: Sec.~\ref{sec:results} shows it is
statistically indistinguishable from the stock recipe.

\subsection{Experimental design}

The plan is declared in a single machine-readable file and executed by a
resumable runner that builds each derived dataset on demand, trains, evaluates
on the held-out split, and appends to a summary table. Configurations are hashed
after defaults are applied, so a configuration shared between factors is trained
once and reused. Every number reported here is regenerated from the run
artefacts by a single script.

Training is deterministic: two independent trainings of one configuration, in
separate processes thirty hours apart, produced per-frame Dice identical to
machine precision. The variance reported in Sec.~\ref{sec:noise} is therefore
attributable to seed choice alone, not to nondeterministic kernels or data
ordering.
